\documentclass[11pt,a4paper]{article}

% ── packages ──────────────────────────────────────────────────────────────
\usepackage[a4paper, margin=2.2cm]{geometry}
\usepackage{amsmath, amssymb}
\usepackage{graphicx}
\usepackage{booktabs}
\usepackage{hyperref}
\usepackage{xcolor}
\usepackage{caption}
\usepackage{subcaption}
\usepackage{parskip}
\usepackage{microtype}
\usepackage{authblk}
\usepackage{titlesec}
\usepackage{enumitem}
\usepackage{float}

\hypersetup{
  colorlinks=true,
  linkcolor=blue!60!black,
  citecolor=blue!60!black,
  urlcolor=blue!60!black
}

\captionsetup{font=small, labelfont=bf, width=0.92\linewidth}
\setlength{\abovecaptionskip}{4pt}
\setlength{\belowcaptionskip}{2pt}

\titleformat{\section}{\large\bfseries}{}{0em}{}[\titlerule]
\titleformat{\subsection}{\normalsize\bfseries}{}{0em}{}

% ── document ──────────────────────────────────────────────────────────────
\begin{document}

% ── title block ───────────────────────────────────────────────────────────
\begin{center}
  {\LARGE\bfseries Computational Reproduction Report}\\[6pt]
  {\large Collective Motion and Inertial Prey--Predator Dynamics}\\[10pt]
  {\normalsize
    Sneha\\
    Indian Institute of Science Education and Research Kolkata\\[4pt]
    Supervisor: Dr.\ Rumi De, IISER Kolkata\\[6pt]
    \textit{June 2026}
  }
\end{center}

\vspace{4pt}
\hrule
\vspace{10pt}

% ── abstract ──────────────────────────────────────────────────────────────
\begin{abstract}
\noindent
This report presents a computational reproduction of two theoretical studies
on collective animal motion and prey--predator dynamics from the group of
Dr.\ Rumi De at IISER Kolkata. The first paper (De and Chakraborty, 2022)
reviews self-propelled particle models for collective motion in fish schools,
bird flocks, and prey swarms. The second (Chakraborty, Laha, and De, 2022)
introduces inertial forces into the prey--predator model and identifies an
optimal predator-to-prey mass ratio governing the arms race between pursuit
and evasion. All simulations are implemented in Julia (physics kernels) with
Python (NumPy/Matplotlib) for figure generation. Key results reproduced
include the Vicsek disorder-to-order phase transition, the metric versus
topological flocking crossover, the interaction-range dependence of prey
survival, and three distinct survival regimes as a function of predator mass.
\end{abstract}

\vspace{6pt}
\hrule
\vspace{10pt}

% ── 1. introduction ───────────────────────────────────────────────────────
\section{Introduction and Motivation}

Collective motion is ubiquitous in biology: starling murmurations,
sardine bait balls, insect swarms, and migrating cell sheets all exhibit
large-scale coherent behaviour that emerges purely from local interactions
among individuals with no global coordinator. Understanding the rules that
generate this order is a central problem in active matter physics.

This reproduction study covers two papers from Dr.\ Rumi De's group at
IISER Kolkata, both of which use self-propelled particle (SPP) models:

\begin{enumerate}[leftmargin=*, itemsep=2pt]
  \item \textbf{De and Chakraborty (2022)} -- a pedagogical review of SPP
    models for fish schooling, bird flocking, and prey--predator swarming,
    with original results on metric versus topological interactions and the
    role of cooperative interaction range in prey survival.
    \textit{J.\ Biosci.\ \textbf{47}:48.}

  \item \textbf{Chakraborty, Laha, and De (2022)} -- introduces inertial
    forces ($M_{pr}$, $M_{pd}$) into the prey--predator model and shows
    that the predator-to-prey mass ratio $M_{pd}/M_{pr}$ governs three
    distinct regimes: full capture, competitive arms race, and prey escape.
    \textit{arXiv:2208.12280.}
\end{enumerate}

% ── 2. models ─────────────────────────────────────────────────────────────
\section{Theoretical Models}

\subsection{2.1 Vicsek Model (1995)}

The canonical SPP model places $N$ point particles in a 2D periodic box of
side $L$. Each particle moves at fixed speed $v_0$ and at each timestep aligns
its heading with the average heading of all neighbours within radius $R$,
perturbed by noise $\eta$:
\begin{align}
  \mathbf{x}_i(t+\Delta t) &= \mathbf{x}_i(t) + \mathbf{v}_i(t)\,\Delta t
    \pmod{L}, \\
  \theta_i(t+\Delta t) &= \langle\theta\rangle_{R} + \Delta\theta,
    \quad \Delta\theta \sim \mathrm{Uniform}\!\left(-\tfrac{\eta}{2},\tfrac{\eta}{2}\right).
\end{align}
The order parameter $\varphi = N^{-1}|\sum_i e^{i\theta_i}|$ measures global
alignment. A continuous phase transition from $\varphi\approx 0$ (disordered)
to $\varphi\approx 1$ (ordered) occurs as $\eta$ decreases or density $\rho$
increases.

\subsection{2.2 Metric and Topological Flocking (Kumar and De, 2021)}

The continuous-time generalisation introduced by Kumar and De (2021) writes:
\begin{equation}
  \frac{d\mathbf{v}_i}{dt}
  = \frac{\alpha}{N_\mathrm{in}}\sum_{N_\mathrm{in}}(\mathbf{v}_j - \mathbf{v}_i)
    - \gamma\,\mathbf{v}_i + \boldsymbol{\xi},
\end{equation}
where $\alpha$ is the alignment strength, $\gamma$ the viscous drag, and
$\boldsymbol{\xi}$ Gaussian noise. In the \textit{metric} rule $N_\mathrm{in}$
counts all neighbours within radius $R$; in the \textit{topological} rule
$N_\mathrm{in} = N_r$ fixed nearest neighbours. Topological interactions
require a smaller threshold $N_r$ to achieve the same ordered state, and the
advantage becomes more pronounced at small flock size.

\subsection{2.3 Overdamped Prey--Predator Swarm (Chakraborty et al., 2020)}

In the overdamped (zero-inertia) limit with friction $\mu=1$:
\begin{align}
  \frac{d\mathbf{r}_i}{dt}
    &= \mathbf{f}_i^{\mathrm{pp}} + \mathbf{f}_i^{\mathrm{pred}}, \\
  \mathbf{f}_i^{\mathrm{pp}}
    &= \frac{1}{N_\mathrm{in}}\sum_{j\in R_\mathrm{int}}
      \!\left[\beta(\mathbf{r}_j - \mathbf{r}_i)
      - \frac{\alpha(\mathbf{r}_j-\mathbf{r}_i)}{|\mathbf{r}_j-\mathbf{r}_i|^2}
      \right], \\
  \mathbf{f}_i^{\mathrm{pred}}
    &= -\frac{\gamma(\mathbf{r}_p - \mathbf{r}_i)}{|\mathbf{r}_p-\mathbf{r}_i|^2},
    \qquad
  \frac{d\mathbf{r}_p}{dt}
    = \frac{\delta}{N}\sum_i
      \frac{\mathbf{r}_i - \mathbf{r}_p}{|\mathbf{r}_i-\mathbf{r}_p|^3}.
\end{align}
The interaction range $R_\mathrm{int}$ controls which escape pattern emerges.
At $R_\mathrm{int}=0$ prey scatter; at intermediate $R_\mathrm{int}$ a ring
forms around the predator; at large $R_\mathrm{int}$ the cohesive group
is tracked and captured.

\subsection{2.4 Inertial Prey--Predator Model (Chakraborty, Laha, and De, 2022)}

Reintroducing inertia through scaled masses $M_{pr}$ and $M_{pd}$:
\begin{align}
  M_{pr}\frac{d^2\mathbf{R}_i}{dT^2}
    &= -\frac{d\mathbf{R}_i}{dT}
      + \frac{1}{N_\mathrm{sur}}\sum_j
        \!\left[\frac{\alpha'\,(\mathbf{R}_i-\mathbf{R}_j)}{|\mathbf{R}_i-\mathbf{R}_j|^2}
        - \beta'(\mathbf{R}_i-\mathbf{R}_j)\right]
      + \frac{\gamma'(\mathbf{R}_i-\mathbf{R}_p)}{|\mathbf{R}_i-\mathbf{R}_p|^2},
    \label{eq:prey-inertial}\\
  M_{pd}\frac{d^2\mathbf{R}_p}{dT^2}
    &= -\frac{d\mathbf{R}_p}{dT}
      - \frac{\delta'}{N_\mathrm{sur}}\sum_i
        \frac{\mathbf{R}_p - \mathbf{R}_i}{|\mathbf{R}_p-\mathbf{R}_i|^3}.
    \label{eq:pred-inertial}
\end{align}
Fixed parameters: $\alpha'=\beta'=1$, $\gamma'=0.2$, $R_\mathrm{kill}=0.01$.
The light predator ($M_{pd}\ll 1$) turns rapidly and captures all prey;
the heavy predator ($M_{pd}\gg 1$) cannot manoeuvre and the prey escape via
the F-maneuver (fountain effect). An intermediate competitive regime exists
where some prey survive and some do not.

% ── 3. numerical implementation ───────────────────────────────────────────
\section{Numerical Implementation}

All physics kernels are written in \textbf{Julia} (v1.9+) using fully
vectorised NumPy-equivalent array operations for speed. Figure generation
uses \textbf{Python} (NumPy, Matplotlib). No per-particle Python loops appear
in any simulation inner loop; all pairwise force matrices are computed via
broadcasting.

\begin{table}[H]
\centering
\caption{Simulation parameters used throughout this reproduction.}
\begin{tabular}{llll}
\toprule
\textbf{Model} & \textbf{Parameter} & \textbf{Value} & \textbf{Notes} \\
\midrule
Vicsek & $N$ & 300 & particles \\
       & $v_0$ & 0.03 & speed \\
       & $R$ & 1.0 & interaction radius \\
       & $\eta$ & 0.1 or 2.0 & low/high noise \\
       & $L$ & 5, 7, 25 & box size \\
\midrule
Flocking & $\alpha$ & 1.5 & alignment strength \\
(Kumar \& De) & $\gamma$ & 0.5 & drag \\
       & $N$ & 50, 100, 200 & flock sizes compared \\
\midrule
Overdamped & $N$ & 200 & prey \\
swarm & $\beta$ & 0.8 & prey-prey attraction \\
       & $\alpha$ & 0.5 & prey-prey repulsion \\
       & $\gamma$ & 0.1 & prey-predator repulsion \\
       & $\delta$ & 0.5 & predator strength \\
       & $R_\mathrm{int}$ & 0 to 2.0 & swept \\
\midrule
Inertial & $N$ & 200 & prey \\
model & $\alpha'=\beta'$ & 1.0 & interaction strengths \\
       & $\gamma'$ & 0.2 & prey-predator coupling \\
       & $\delta'$ & 0.4 (weak) / 2.5 (strong) & predator strength \\
       & $M_{pd}$ & 0.1 to 100 & swept \\
       & $dt$ & 0.01--0.02 & timestep \\
\bottomrule
\end{tabular}
\end{table}

% ── 4. results: paper 1 ───────────────────────────────────────────────────
\section{Results: De and Chakraborty (2022)}

\subsection{4.1 Vicsek Phase Transition}

Figure~\ref{fig:phi_eta} shows the order parameter $\varphi$ as a function
of noise $\eta$ for $N=300$, $L=7$, $R=1$, reproducing the continuous
disorder-to-order transition. The transition occurs near $\eta\approx 0.8$
for these parameters, consistent with the original Vicsek et al.\ (1995)
result.

\begin{figure}[H]
  \centering
  \includegraphics[width=0.55\linewidth]{figs/static1_phi_vs_eta.png}
  \caption{Order parameter $\varphi$ vs noise strength $\eta$. As $\eta$
    decreases (right to left), the system transitions from a disordered phase
    ($\varphi\approx 0$) to a collectively ordered flock ($\varphi\approx 1$).
    Red dashed line marks $\varphi=0.5$.}
  \label{fig:phi_eta}
\end{figure}

\subsection{4.2 Metric vs Topological Flocking}

Figure~\ref{fig:metric_topo} reproduces Fig.\ 3 of De and Chakraborty (2022).
For metric interactions (left panel), a threshold radius $R^*$ is required
for order, and $R^*$ decreases with increasing $N$ as each particle gains more
neighbours. For topological interactions (right panel), the threshold number
of neighbours $N_r^*$ is independent of flock size at lower speeds. In both
cases the ordered state is reached more easily with larger groups.

\begin{figure}[H]
  \centering
  \includegraphics[width=0.92\linewidth]{figs/static2_metric_vs_topo.png}
  \caption{Order parameter $\varphi$ vs (a) metric interaction radius $R$
    and (b) topological neighbour count $N_r$, for flock sizes $N=50, 100,
    200$. Topological interactions achieve the same ordered state with fewer
    effective neighbours, and the threshold is size-independent at low speed.}
  \label{fig:metric_topo}
\end{figure}

\subsection{4.3 Overdamped Prey Swarm: Escape Patterns}

Figure~\ref{fig:swarm_patterns} shows four qualitatively distinct escape
patterns that emerge as $R_\mathrm{int}$ increases. At $R_\mathrm{int}=0$
prey scatter incoherently. At $R_\mathrm{int}=0.5$ prey organise into
a ring encircling the predator, confusing its direction of attack. At
$R_\mathrm{int}=1.2$ the group splits into two coherent subgroups flanking
the predator. At $R_\mathrm{int}=2.0$ the cohesive group is easily tracked
and chased.

\begin{figure}[H]
  \centering
  \includegraphics[width=0.85\linewidth]{figs/snap1_swarm_Rint_sweep.png}
  \caption{Escape patterns for $N=200$ prey at increasing interaction radius
    $R_\mathrm{int}$ (rows). Columns show the initial state ($T=0$),
    an intermediate time, and the established pattern. Blue dots: alive
    prey; grey dots: killed prey; red circle: predator. Dashed circle
    indicates $R_\mathrm{int}$.}
  \label{fig:swarm_patterns}
\end{figure}

\subsection{4.4 Survival vs Interaction Range}

Figure~\ref{fig:survival_rint} reproduces Fig.\ 6 of De and Chakraborty
(2022). Prey survival $N_\mathrm{sur}$ is maximised at intermediate
$R_\mathrm{int} \approx 0.5$--$1.3$ across all predator strengths tested.
Too little cooperation leaves prey unable to coordinate; too much cooperation
creates a cohesive target that the predator tracks efficiently.

\begin{figure}[H]
  \centering
  \includegraphics[width=0.60\linewidth]{figs/static3_survival_vs_rint.png}
  \caption{Number of surviving prey $N_\mathrm{sur}$ as a function of
    interaction radius $R_\mathrm{int}$, for three predator strengths
    $\delta$. The green band marks the optimal cooperation range. $N=150$,
    averaged over 3 realisations.}
  \label{fig:survival_rint}
\end{figure}

% ── 5. results: paper 2 ───────────────────────────────────────────────────
\section{Results: Chakraborty, Laha, and De (2022)}

\subsection{5.1 Weak Predator: Ring Formation and Chasing Transition}

Figure~\ref{fig:weak_pred} shows how predator mass $M_{pd}$ controls whether
a stable prey ring forms around the predator ($\delta'=0.4$). For
$M_{pd}=2.0$ the ring is fully established by $T=150$ and persists. As
$M_{pd}$ increases to 3.0 and beyond, the ring destabilises: higher inertia
slows the predator's turning response, breaking the rotational symmetry of
the ring configuration. By $M_{pd}=5.0$ only active chasing occurs from the
outset. This matches the stability analysis in the paper, which shows an
eigenvalue crossing zero at $M_{pd}\approx 3.0$ (for $M_{pr}=1$).

\begin{figure}[H]
  \centering
  \includegraphics[width=0.88\linewidth]{figs/snap2_inertial_weak_Mpd_sweep.png}
  \caption{Weak predator ($\delta'=0.4$, $M_{pr}=1$): effect of predator mass
    on escape pattern. Each row shows a different $M_{pd}$ value; columns are
    the initial state, an intermediate time (ring/pattern forming), and the
    established pattern. Arrows show prey velocity directions. The stable ring
    at $M_{pd}=2.0$ (top row, $T=150$) is the key result.}
  \label{fig:weak_pred}
\end{figure}

\subsection{5.2 Strong Predator: Escape Trajectories}

Figure~\ref{fig:strong_pred} shows the three qualitatively different escape
trajectories that emerge for the strong predator ($\delta'=2.5$) as $M_{pd}$
is varied. The light predator ($M_{pd}=0.1$) turns rapidly and maintains
close pursuit inside the herd. The medium predator ($M_{pd}=1.0$) causes the
prey group to split and merge. The heavy predator ($M_{pd}=100$) overshoots
(visible from the orange trajectory trail) and cannot reorient fast enough
to prevent the F-maneuver, whereby prey arc around and escape.

\begin{figure}[H]
  \centering
  \includegraphics[width=0.88\linewidth]{figs/snap3_inertial_strong_Mpd_sweep.png}
  \caption{Strong predator ($\delta'=2.5$, $M_{pr}=1$): escape trajectories
    for four predator masses. Orange trail shows the predator's path. The
    middle column shows the predator inside the prey herd. The right column
    shows the established escape pattern: capture (top), splitting (second),
    arc escape (third), and F-maneuver with full prey escape (bottom).}
  \label{fig:strong_pred}
\end{figure}

\subsection{5.3 Order Parameter $\varphi(T)$}

The order parameter $\varphi(T) = N_\mathrm{sur}^{-1}\sum_i \cos\theta_i$
measures the mean cosine of the angle between prey velocity $i$ and the
predator velocity. Figure~\ref{fig:phi_T} reproduces Fig.\ 5 of the paper.
For $M_{pd}=0.1$ (light predator) $\varphi$ oscillates rapidly as the
predator continuously reverses direction to chase. For $M_{pd}=100$
(heavy predator) $\varphi$ drops to $-1$ and stays there: prey and predator
move in opposite directions for an extended period, indicating successful
escape.

\begin{figure}[H]
  \centering
  \includegraphics[width=0.85\linewidth]{figs/static5_phi_timeseries.png}
  \caption{Order parameter $\varphi(T)$ for four predator masses.
    $\varphi=+1$: prey and predator move in the same direction.
    $\varphi=-1$: prey and predator move in opposite directions (escape).
    Green/red shading indicates cooperative vs opposing motion.}
  \label{fig:phi_T}
\end{figure}

\subsection{5.4 Survival Phase Diagram}

Figure~\ref{fig:survival_mass} reproduces Fig.\ 6 of Chakraborty, Laha, and
De (2022). Panel (a) shows $N_\mathrm{sur}$ vs $M_{pd}$ for four prey masses.
Lighter prey survive better at low $M_{pd}$ because their own agility
outcompetes the light predator. Panel (b) collapses the results onto the mass
ratio $M_{pd}/M_{pr}$, revealing three universal regimes separated by order-of-magnitude boundaries:

\begin{itemize}[itemsep=2pt]
  \item \textbf{Killed regime} ($M_{pd}/M_{pr} \lesssim 0.4$): the agile
    predator captures all prey.
  \item \textbf{Competitive regime} ($0.4 \lesssim M_{pd}/M_{pr} \lesssim 10$):
    some prey survive, some do not -- the arms race.
  \item \textbf{Survival regime} ($M_{pd}/M_{pr} \gtrsim 10$): the sluggish
    predator cannot manoeuvre and all prey escape.
\end{itemize}

\begin{figure}[H]
  \centering
  \includegraphics[width=0.92\linewidth]{figs/static4_survival_vs_mass.png}
  \caption{(a) $N_\mathrm{sur}$ vs predator mass $M_{pd}$ for four prey
    masses $M_{pr}$, $\delta'=2.5$, $N=80$. (b) $N_\mathrm{sur}$ vs the
    mass ratio $M_{pd}/M_{pr}$, collapsing onto three universal regimes.}
  \label{fig:survival_mass}
\end{figure}

% ── 6. discussion ─────────────────────────────────────────────────────────
\section{Discussion}

\textbf{Vicsek transition.} The reproduction confirms the first-order-like
nature of the transition: $\varphi$ rises steeply over a narrow range of
$\eta$, consistent with the finite-size scaling analysis of the original paper.

\textbf{Topological advantage.} The topological rule is more robust than
the metric rule at small $N$: removing metric distance as the interaction
criterion means that even sparse flocks can coordinate, matching the
empirical finding by Ballerini et al.\ (2008) that starlings interact with
approximately 6--7 nearest neighbours regardless of separation.

\textbf{Optimal cooperation range.} The non-monotonic dependence of
$N_\mathrm{sur}$ on $R_\mathrm{int}$ is a non-trivial result. It reflects
a tradeoff: enough cooperation is needed to produce coordinated evasion
(ring or splitting), but too much cooperation makes the group a coherent,
easily-tracked target. The optimal $R_\mathrm{int}$ shifts slightly with
predator strength $\delta$, consistent with the paper's Fig.\ 6.

\textbf{Inertial arms race.} The identification of three regimes on the
$M_{pd}/M_{pr}$ axis is the central result of Chakraborty, Laha, and De
(2022). The competitive intermediate regime -- where field studies observe
cheetah-impala and lion-zebra interactions -- corresponds to $M_{pd}/M_{pr}
\sim 1$--$10$, consistent with empirically measured body-mass ratios for
those predator-prey pairs (Wilson et al., 2018).

\textbf{Numerical considerations.} The overdamped model is sensitive to
timestep: $dt=0.005$ was required for stable integration. The inertial model
required $dt\leq 0.02$ to avoid velocity blow-up at low $M_{pr}$. Vectorised
pairwise force computation scales as $O(N^2)$; for $N=200$ this runs in
under 10 ms per timestep in Julia.

% ── 7. conclusion ─────────────────────────────────────────────────────────
\section{Conclusion}

All key results from both papers have been successfully reproduced:
the Vicsek phase transition, metric/topological flocking crossover,
the four overdamped escape patterns and their survival statistics, and
the three inertial survival regimes. The Julia source modules
(\texttt{vicsek\_model.jl}, \texttt{flocking\_interactions.jl},
\texttt{prey\_predator\_swarm.jl}, \texttt{inertial\_prey\_predator.jl})
and Python animation scripts are available in the accompanying GitHub
repository.

Future extensions include: (i) implementing finite-size scaling for the
Vicsek transition, (ii) incorporating ABP (active Brownian particle)
dynamics for the prey in porous media as studied in related work, and
(iii) extending to 3D geometries relevant for fish schools.

% ── references ────────────────────────────────────────────────────────────
\section*{References}

\begin{enumerate}[label={[\arabic*]}, leftmargin=*, itemsep=3pt]

\item De R and Chakraborty D (2022). Collective motion: influence of local
  behavioural interactions among individuals.
  \textit{J.\ Biosci.}\ \textbf{47}:48.

\item Chakraborty D, Laha A, and De R (2022). Effect of inertia on the
  evasion and pursuit dynamics of prey swarms and the emergence of an
  optimal mass ratio for the predator-prey arms race.
  \textit{arXiv:}2208.12280.

\item Vicsek T, Czir\'{o}k A, Ben-Jacob E, Cohen I, and Shochet O (1995).
  Novel type of phase transition in a system of self-driven particles.
  \textit{Phys.\ Rev.\ Lett.}\ \textbf{75}:1226.

\item Kumar V and De R (2021). Efficient flocking: metric versus topological
  interactions.
  \textit{R.\ Soc.\ Open Sci.}\ \textbf{8}:58.

\item Chakraborty D, Bhunia S, and De R (2020). Survival chances of a prey
  swarm: how the cooperative interaction range affects the outcome.
  \textit{Sci.\ Rep.}\ \textbf{10}:8362.

\item Ballerini M et al.\ (2008). Interaction ruling animal collective
  behavior depends on topological rather than metric distance.
  \textit{Proc.\ Natl Acad.\ Sci.\ USA}\ \textbf{105}:1232.

\item Wilson AM et al.\ (2018). Biomechanics of predator-prey arms race in
  lion, zebra, cheetah and impala.
  \textit{Nature}\ \textbf{554}:183.

\end{enumerate}

\end{document}
